\section{Introduction}

% ------------------------------------------------------------
\begin{frame}
  \frametitle{About the Lecturer}

  \begin{columns}
    \column{0.6\textwidth}
      \begin{block}{Mathis Neunzig}
        \begin{itemize}
          \item B.Sc.\ Business Information Technology\\
                (DHBW Mannheim)
          \item ABAP Developer / Trainer /\\
                Student-Recruiting @ SAP
          \item Also teaches Software Engineering\\
                at DHBW
          \item Based in Heidelberg
        \end{itemize}
      \end{block}

    \column{0.38\textwidth}
      \begin{block}{Contact}
        \footnotesize
        \textbf{Email}\\
        mathis.neunzig@gmail.com\\[6pt]
        \textbf{Phone}\\
        +49\,174\,9885992
      \end{block}
  \end{columns}

  \vspace{0.5em}
  Your lecturer combines hands-on industry experience at one of the world's largest software companies with a genuine passion for teaching. Feel free to reach out at any time -- whether for course-related questions or simply to connect professionally. Getting to know your students is an important part of good teaching.
\end{frame}

% ------------------------------------------------------------
\begin{frame}
  \frametitle{Getting to Know You}

  \begin{block}{Let's go around the room -- please tell us:}
    \begin{enumerate}
      \item Your \textbf{name}
      \item Your \textbf{employer} (dual-study company)
      \item Your \textbf{favourite animal}
      \item Your \textbf{favourite colour}
      \item Your \textbf{favourite song or album} right now
      \item Your \textbf{favourite programming language} (so far)
    \end{enumerate}
  \end{block}

  \vspace{0.5em}
  These introductions serve a dual purpose: they help break the ice at the start of a new semester and give the lecturer a sense of the group's technical background and personality. The answers to the last question in particular can be very revealing -- knowing which languages you already feel comfortable with will help contextualise the SQL content that follows later in the course. There are no wrong answers here, so be honest and have fun with it!
\end{frame}

% ------------------------------------------------------------
\begin{frame}
  \frametitle{What Will We Cover?}

  This course is structured around three interconnected topic groups that build on each other progressively throughout the semester.

  \begin{block}{1 -- Databases (Foundations)}
    \begin{itemize}
      \item What is a database / DBMS?
      \item History of database systems
      \item Relational algebra
      \item Relational databases
      \item Data modelling \& the DB design process
    \end{itemize}
  \end{block}

  \begin{block}{2 -- SQL}
    \begin{itemize}
      \item Query languages overview
      \item DQL, DML, DDL, DCL
      \item Extensive hands-on exercises
    \end{itemize}
  \end{block}

  \begin{block}{3 -- Advanced Topics}
    \begin{itemize}
      \item Transactions \& multi-user access
      \item Backup \& Recovery
      \item Security, ORM, Spring Data JPA
    \end{itemize}
  \end{block}
\end{frame}

% ------------------------------------------------------------
\begin{frame}
  \frametitle{Semester Schedule}

  \scriptsize
  \begin{tabularx}{\linewidth}{@{}l l l X@{}}
    \toprule
    \textbf{Session} & \textbf{Date} & \textbf{Time} & \textbf{Topics} \\
    \midrule
    S3-1  & 09.09 & 12:30--15:00 & Introduction; Course overview; Exam info; DB Management System; Data Models \\
    S3-2  & 15.09 & 14:00--17:15 & Relational Databases; Basic Relational Algebra \\
    S3-3  & 19.09 & 09:00--12:15 & Advanced Relational Algebra; Normal Forms (1NF, 2NF, 3NF) \\
    S3-4  & 26.09 & 09:00--12:15 & Normal Forms (BCNF, 4NF, 5NF, 6NF) \\
    S3-5  & 29.09 & 14:00--17:15 & ERM Design \\
    S3-6  & 13.10 & 14:00--17:15 & SQL (DML) \\
    S3-7  & 20.10 & 14:00--17:15 & SQL (DML) \\
    S3-8  & 27.10 & 14:00--17:15 & SQL (DCL); CAP-Theorem \\
    S3-9  & 03.11 & 14:00--17:15 & Backup \& Recovery; Performance; Concurrency \\
    S3-10 & 10.11 & 14:00--17:15 & Security; Object-Relational Mapping \\
    S3-11 & 17.11 & 14:00--17:15 & Exam preparation \\
    S3-12 & 26.11 & 10:00--11:00 & \textbf{Written Exam} \\
    \bottomrule
  \end{tabularx}

  \vspace{0.4em}
  \normalsize
  The schedule is designed to move from abstract theory (data models, relational algebra, normal forms) towards increasingly practical content (SQL, ORM, security). Attendance is strongly recommended throughout -- later sessions build directly on concepts introduced in earlier ones.
\end{frame}

% ------------------------------------------------------------
\begin{frame}
  \frametitle{Organization}

  \begin{block}{Class Representative}
    Please elect a class representative by the end of the first session. The representative acts as the primary point of contact between the student group and the lecturer for administrative matters, schedule changes, and feedback.
  \end{block}

  \begin{block}{Mailing List \& Contact}
    \begin{itemize}
      \item All course announcements will be sent via the mailing list -- make sure your email address is registered.
      \item Lecturer contact: \texttt{mathis.neunzig@gmail.com} \quad +49\,174\,9885992
    \end{itemize}
  \end{block}

  \begin{alertblock}{Emergency Contact}
    If you have an urgent issue that cannot wait for the next session -- for example a technical problem on the day of an assignment submission -- please use the phone number above. For non-urgent matters, email is always preferred.
  \end{alertblock}

  Good organisation at the start of a course saves a great deal of confusion later. Make sure you know who your class representative is and that the mailing list is working correctly for you from day one.
\end{frame}

% ------------------------------------------------------------
\begin{frame}
  \frametitle{Rules}

  \begin{block}{General Conduct}
    Standard DHBW rules of conduct apply throughout all sessions.
  \end{block}

  \begin{itemize}
    \item \textbf{Absences:} Please give advance notice if you cannot attend. A quick message the day before is always appreciated -- it helps the lecturer plan group exercises appropriately.
    \item \textbf{Deadlines:} If you foresee that you cannot meet an assignment deadline, contact the lecturer \emph{before} the deadline, not after. Extensions can often be arranged if asked in advance.
    \item \textbf{Open door:} Feel free to reach out at any time with questions, concerns, or feedback -- about the content, the pace, or anything else. There are no silly questions in this course.
    \item \textbf{Group exercises:} Active participation during group exercises is expected and contributes to your own learning as well as the group's. Sitting passively during a group task is not acceptable.
  \end{itemize}

  These rules are based on mutual respect. The expectation is not perfection but honest communication -- if something is not working for you, say so early rather than letting it become a bigger problem.
\end{frame}

% ------------------------------------------------------------
\begin{frame}
  \frametitle{Assessment: Written Exam}

  \begin{block}{Key Facts}
    \begin{itemize}
      \item \textbf{Weight:} 40 points \quad(40\,\% of final grade)
      \item \textbf{Timing:} End of 3rd semester (S3-12, 26.11)
      \item \textbf{Duration:} 60 minutes
    \end{itemize}
  \end{block}

  \begin{alertblock}{Exam Scope}
    Topics marked with \textbf{(!)} in the slides will \emph{not} appear in the written exam. These are advanced or supplementary topics included for completeness and professional context, but you will not be tested on them under exam conditions.
  \end{alertblock}

  \begin{examples}{Permitted Aids}
    You may bring \textbf{one handwritten A4 sheet} (both sides allowed) as a personal cheat sheet. It must be handwritten -- printed or typed sheets are not permitted. Writing your own cheat sheet is itself an effective revision strategy.
  \end{examples}

  The written exam tests your understanding of foundational database concepts, relational algebra, normalisation theory, and SQL. Focus on understanding the \emph{why} behind each concept, not just memorising syntax -- many exam questions require you to reason through a scenario.
\end{frame}

% ------------------------------------------------------------
\begin{frame}
  \frametitle{Assessment: Case Studies / Assignments}

  \begin{block}{Overview -- 60 points total (60\,\% of final grade)}
    \begin{tabularx}{\linewidth}{@{}l X r@{}}
      \toprule
      \textbf{Assignment} & \textbf{Description} & \textbf{Points} \\
      \midrule
      \multicolumn{3}{@{}l}{\textit{Semester 3}} \\
      Assignment 1 & Individual work            & 10 pts \\
      Assignment 2 & Groups of 2--3 students    & 10 pts \\
      \multicolumn{3}{@{}l}{\textit{Semester 4}} \\
      Assignment 3 & Individual work            & 10 pts \\
      Assignment 4 & Final project (groups 2--3) & 30 pts \\
      \bottomrule
    \end{tabularx}
  \end{block}

  \begin{block}{Submission Policy (Assignments 1--3)}
    \begin{itemize}
      \item \textbf{1st submission} receives a grade \emph{and} written feedback.
      \item \textbf{Optional 2nd submission} is possible after reviewing the feedback.
      \item The 2nd submission grade replaces the 1st -- use this opportunity wisely.
    \end{itemize}
  \end{block}

  The two-submission model is designed to support your learning: you are encouraged to engage with the feedback and genuinely improve your solution, not just tweak minor details. Assignment 4 is the most substantial deliverable and will give you the opportunity to apply everything you have learned in a realistic scenario.
\end{frame}

% ------------------------------------------------------------
\begin{frame}
  \frametitle{Literature}

  \begin{block}{No textbook purchase is required to pass this course.}
    All essential material is covered in the lecture slides and exercises.
  \end{block}

  \textbf{Recommended reading (for those who want to go deeper):}

  \begin{itemize}
    \item \textbf{[DE]} Kemper \& Eickler -- \textit{Datenbanksysteme}\\
          \footnotesize The German-language standard reference for database courses at German universities. Thorough and mathematically precise.\normalsize

    \item Garcia-Molina, Ullman \& Widom -- \textit{Database Systems: The Complete Book}, 2nd~ed.\\
          \footnotesize Comprehensive English-language textbook covering theory and practice in depth.\normalsize

    \item Foster \& Godbole -- \textit{Database Systems: A Pragmatic Approach}, 2nd~ed.\\
          \footnotesize A more application-oriented introduction, well-suited as a companion to this course.\normalsize

    \item \textbf{[DE]} Gumm \& Sommer -- \textit{Einführung in die Informatik}, 10th~ed.\\
          \footnotesize Broader CS introduction, useful for contextualising databases within the wider field.\normalsize
  \end{itemize}

  Reading around the subject -- even dipping into one or two chapters -- significantly deepens your understanding and will help you engage more confidently with the material during lectures.
\end{frame}
